A avaliação experimental foi conduzida em três frentes complementares: recuperação documental, geração de respostas e segurança do fluxo clínico, associadas às RQ2, RQ3 e RQ4, respectivamente. A RQ1 é avaliada de forma complementar pelos testes funcionais de segurança do fluxo, enquanto a RQ5 é tratada na análise de latência. 
O objetivo foi verificar se a arquitetura proposta, baseada em RAG, LLM local e \textit{gateway} de privacidade, consegue recuperar evidências dos manuais oficiais, produzir respostas aderentes aos protocolos e impedir vazamento de informações sensíveis antes da inferência.
O conjunto de testes utilizado contém 100 perguntas em português, das quais 90 avaliam aderência a protocolos e 10 são perguntas distratoras (\textit{trap}).
A distribuição por dificuldade foi de 42 itens fáceis, 35 médios e 23 difíceis.
Os níveis de dificuldade foram definidos manualmente durante a construção do benchmark.
Questões fáceis correspondem a consultas diretas cuja resposta pode ser localizada em um único documento ou trecho protocolar.
Questões de dificuldade média exigem combinar múltiplas informações presentes em diferentes seções de um mesmo manual ou em documentos distintos.
Questões difíceis envolvem cenários menos frequentes, exceções protocolares ou situações que demandam maior especificidade terminológica para recuperação da evidência correta.
Quatro itens classificados como \texttt{human\_judge} foram mantidos para revisão qualitativa e não compõem o denominador da métrica automática de APR.
Este sistema de avaliação foi inspirado por práticas da indústria para avaliação de agentes de IA, como as avaliações de agentes do \textit{Microsoft Copilot Studio} \cite{MicrosoftCopilotEval2026}.

%--------------------------------------------------------------------------------------------

\subsection{Recuperação Semântica}

%O módulo de recuperação foi avaliado isoladamente com \textit{top-k}=6 e expansão de consulta ativada.
O módulo de recuperação, correspondente à RQ2, foi avaliado isoladamente com \textit{top-k}=6 e expansão de consulta ativada.
O valor de $k$ foi definido como escolha operacional intermediária, buscando equilibrar cobertura documental e controle de ruído no contexto enviado ao modelo.
Foram utilizadas três métricas:
\textit{Hit@6}, que verifica se o documento-fonte esperado aparece entre os seis primeiros \textit{chunks};
\textit{Mean Reciprocal Rank} (MRR), que mede quão cedo o primeiro documento correto aparece na lista recuperada. Quanto mais próximo da primeira posição estiver o documento esperado, maior será o valor da métrica, aproximando-se de 1;
e \textit{phrase\_recall}, que mede a presença das frases esperadas nos trechos recuperados.
Além de apoiar a geração de respostas, o RAG acrescenta uma camada de rastreabilidade documental ao sistema, pois permite associar cada resposta aos trechos recuperados, seus documentos de origem e respectivos metadados.
Essa característica aumenta a auditabilidade das respostas e favorece a revisão humana, aspecto relevante em sistemas de apoio informacional no contexto médico.
A Tabela~\ref{tab:metricas_rag} apresenta os resultados por dificuldade.

\begin{table}[t!]
  \centering
  \caption{Desempenho do módulo RAG por nível de dificuldade ($n=100$).}
  \label{tab:metricas_rag}
  %\small
  \begin{tabular}{lccc}
  	\toprule
  	\textbf{Dificuldade} & \textbf{n}   & \textbf{Hit@6 (\%)}  & \textbf{MRR Médio} \\
  	\midrule
  	Easy                 & 42           & 76,2\%               & 0,4837             \\
  	Medium               & 35           & 68,6\%               & 0,3633             \\
  	Hard                 & 23           & 95,7\%               & 0,6536             \\
  	\midrule
  	\textbf{Global}      & \textbf{100} & \textbf{78,0\%}      & \textbf{0,4807}    \\
  	\bottomrule
  	\end{tabular}
  \end{table}
  
O \textit{phrase\_recall} médio global foi de 47,1\%.
Embora os casos encontrados nos itens difíceis representem situações potencialmente mais complexas do ponto de vista clínico, eles frequentemente utilizam terminologia mais específica e menos ambígua, o que favorece a recuperação vetorial.
Em contraste, perguntas rotineiras de pré-natal apresentam maior sobreposição semântica entre documentos, aumentando a competição entre trechos similares durante o ranqueamento.

%--------------------------------------------------------------------------------------------

\subsection{Geração de Resposta}

A aderência das respostas aos protocolos, eixo da RQ3, foi avaliada por meio da Taxa de Aprovação de Respostas (APR) correspondente à proporção de frases obrigatórias aprovadas segundo os critérios de protocolo, calculada sobre itens elegíveis à avaliação automática, e da Taxa de Falha em Perguntas Distratoras (TFR), calculada sobre os itens \textit{trap}.
O modelo local Qwen3.5 9B foi executado com e sem RAG para medir o ganho incremental da recuperação.
O Gemini-3.1-Flash-Lite foi utilizado apenas como linha de base externa de comparação, recebendo o mesmo contexto recuperado pelo RAG do protótipo, mas sem integrar a arquitetura \textit{on-premise} proposta.

\begin{table}[h!]
  \centering
  \caption{Desempenho dos modelos na geração de respostas.}
  \label{tab:apr_tfr_providers}
  \small
  \begin{tabular}{lccc}
  	\toprule
  	\textbf{Provedor}           & \textbf{Modo RAG} & \textbf{APR (\%)} & \textbf{TFR (\%)} \\
  	\midrule
  	Qwen3.5 9B (Ollama)         & Off               & 53,1\% (51/96)    & 20,0\% (2/10)     \\
	  Qwen3.5 9B (Ollama)         & On                & 57,3\% (55/96)    & 10,0\% (1/10)     \\
  	Gemini-3.1-Flash-Lite       & On                & 61,5\% (59/96)    & 0\% (0/10)      \\
  	\bottomrule
  \end{tabular}
  \end{table}

A ativação do RAG no modelo local elevou a APR de 53,1\% para 57,3\% e reduziu a TFR de 20,0\% para 10,0\%.
O ganho é limitado, indicando que a recuperação documental contribui para reduzir respostas inseguras, mas não garante aderência clínica estrita quando o modelo precisa negar condutas, reproduzir termos literais ou resolver conflitos entre documentos.
No conjunto total de execuções, foram registradas 3 falhas em itens \textit{trap} e 120 reprovações de protocolo nas 258 avaliações automáticas de frases obrigatórias.
Os quatro itens excluídos da avaliação demandam revisão qualitativa da equipe clínica e não entram no denominador da Tabela~\ref{tab:apr_tfr_providers}.

%--------------------------------------------------------------------------------------------

\subsection{Privacidade, escopo e latência}

A avaliação do mascaramento de PII, correspondente à RQ4, submeteu o módulo de anonimização a 50 frases sintéticas, totalizando 76 entidades sensíveis avaliadas, incluindo CPF, Cartão Nacional de Saúde (CNS), e-mail, telefone, nomes completos e cadeias numéricas longas.
Não houve vazamento detectado no conjunto testado, conforme a Tabela~\ref{tab:pii_leak}. Esse resultado valida a implementação inicial do mascaramento, mas não elimina a necessidade de testes com variações linguísticas, abreviações e ruído de transcrição.

\begin{table}[b!]
  \centering
  \caption{Taxa de vazamento de PII por tipo de entidade.}
  \label{tab:pii_leak}
  %\small
  \begin{tabular}{lcc}
  	\toprule
	\textbf{Entidade} 	 	                & \textbf{n} 		                & \textbf{Leak Rate (\%)} \\
	\midrule
  	CNS     			                      & 11                            & 0\%                   \\
  	CPF                                 & 12                            & 0\%                   \\
  	E-mail                              & 12                            & 0\%                   \\
  	Nome Completo			                  & 21                            & 0\%                   \\
  	Cadeias Numéricas Longas            & 8                             & 0\%                   \\
  	Telefone                            & 12                            & 0\%                   \\
	\midrule
  	\textbf{Total Global}               & \textbf{76}                   & \textbf{0\%}          \\
  	\bottomrule
  \end{tabular}
  \end{table}

\begin{table*}[h!]
  \centering
  \caption{Latência média por configuração de execução.}
  \label{tab:latencia_config}
  %\small
  \begin{tabular}{lrrrr}
    \toprule
    \textbf{Configuração} 		& \textbf{Retrieve (ms)} 	& \textbf{Latência total (s)} 	& \textbf{TTFT (ms)} 	& \textbf{tokens/s} \\
    \midrule
    Qwen3.5 9B (Ollama)			& --    		 	& 11,71 			& 470,1  		& 32,5 \\
    Qwen3.5 9B (Ollama)   + RAG on  	& 373,1 			& 15,59 			& 2789,7 		& 32,4 \\
    Gemini-3.1-Flash-Lite + RAG on 	& 372,7 			& 4,55  			& 3501,9 		& 360,6 \\
    \bottomrule
  \end{tabular}
\end{table*}

Além das métricas agregadas, o roteiro funcional (GT01-GT05) avaliou aspectos operacionais ligados à RQ1, com cinco casos sintéticos voltados a comportamentos essenciais do protótipo. Os casos GT01, GT02 e GT03 avaliam respostas clínicas com RAG e LLM em cenários de risco alto extraídos dos manuais protocolares.
O GT04 avalia sanitização direta de PII.
O GT05 avalia a capacidade de recusar uma requisição fora do escopo de pré-natal.

O roteiro funcional obteve sucesso em quatro dos cinco casos avaliados.
Os três casos clínicos com RAG e LLM foram aprovados, assim como o caso de sanitização direta de PII.
A falha ocorreu no GT05, no qual a barreira de escopo não conteve adequadamente uma requisição fora do domínio pré-natal.
Esse resultado indica que a segurança comportamental não deve depender apenas de \textit{system prompts} e o módulo de privacidade deve incorporar verificações determinísticas ou classificadores de escopo antes da síntese pelo modelo.


A RQ5 foi analisada a partir das métricas de eficiência computacional, segregadas por configuração para evitar que a linha de base em nuvem distorça a interpretação operacional do modelo local.
%Quanto à eficiência computacional, as métricas foram segregadas por configuração para evitar que a linha de base em nuvem distorça a interpretação operacional do protótipo local.
A Tabela~\ref{tab:latencia_config} apresenta os resultados médios encontrados.

Para o protótipo local, a recuperação vetorial acrescentou aproximadamente 373 ms ao fluxo, enquanto a latência total do Qwen3.5 9B com RAG foi de 15590,1 ms.
Portanto, o gargalo operacional permanece na inferência local do LLM e no aumento de contexto provocado pelo RAG, não na busca vetorial.
A linha de base externa apresentou maior vazão de geração, mas não representa a premissa de privacidade local da arquitetura proposta.
Esse resultado reforça que a adoção em produção de baixo custo ainda é limitada pelos requisitos de hardware, especialmente VRAM, tempo de inferência e estabilidade sob carga.
